%% kapa: exam=1 round=19 year=2008 period=2 form=A subject=economics subject_ko=경제원론 seq=36 number=36
%% source: Q-Net 감정평가사 기출문제 (cst00309) artlSeq=1001962 file=extracted/1008/1001962/2008년 제19회 감정평가사 1차시험 기출문제/2교시.hwp
%% conversion: kordoc 4.24.51 -> tools/gen1.py (auto); review=proofread-sonnet; flags=none
\begin{question}{36}
정부가 경기활성화를 위하여 오피스텔 임대업자에게 부과하는 세금을 감면하기로 결정하였다고 하자. 오피스텔 임대시장의 변화에 대한 설명으로 적절한 것은? (단, 오피스텔 임대시장의 수요곡선은 우하향하고, 공급곡선은 우상향한다고 가정한다.)
\begin{choices}
  \item 오피스텔 임대료가 하락한다.
  \item 오피스텔 임대시장의 공급곡선이 왼쪽으로 이동한다.
  \item 오피스텔 임대시장의 수요곡선이 오른쪽으로 이동한다.
  \item 조세감면혜택은 오피스텔 임대시장의 공급자에게만 귀착된다.
  \item 조세감면혜택은 오피스텔 임대시장의 수요자에게만 귀착된다.
\end{choices}
\end{question}
\answer{1}
